%%
%% Copyright (c) 2018-2020 Weitian LI <wt@liwt.net>
%% CC BY 4.0 License
%%

% Chinese version
\documentclass[zh]{resume}

% File information shown at the footer of the last page
\fileinfo{%
  \faCopyright{} 2025--2026, 劲舟梁 \hspace{0.5em}
  \creativecommons{by}{4.0} \hspace{0.5em}
  \githublink{dok4everak47}{LLM-Resume-Template} \hspace{0.5em}
  \faEdit{} \today
}

\name{劲舟}{梁}

\keywords{机器学习, 模型训练, 数据处理, 自动驾驶, Agent}

\tagline{求职意向：机器学习工程师}

\photo[shape=circular, position=right]{2.0cm}{../assets/luongchin.png}

\profile{
  \mobile{18176608865}
  \email{dok4ever123@126.com}
  \iconlink{\faGithub}{https://github.com/dok4everak47}{dok4everak47}
}

\begin{document}
\makeheader

%======================================================================
\sectionTitle{个人总结}{\faUser}
%======================================================================
具有数学与计算科学背景的机器学习实践者，聚焦模型训练数据构建、质量评估与工程实现。在自动驾驶 3D 感知数据标注策略、RLHF 偏好数据构建与多模态 AIGC 评估方向有落地经验；从零复现 Claude Code 核心架构，对 LLM Agent 的执行循环、Tool Use 与 Sub-Agent 分治机制有深入理解，能独立完成数据处理流水线搭建与模型评测分析。

%======================================================================
\sectionTitle{专业技能}{\faWrench}
%======================================================================
\begin{itemize}
    \item \textbf{编程语言}: Python, TypeScript, Node.js, Pandas, NumPy, scikit-learn
    \item \textbf{机器学习 / 模型}: 训练数据构建与清洗、模型输出分析与质量评估、RLHF 偏好数据构建
    \item \textbf{数据处理}: 多源日志解析与噪声清洗、标注策略制定、特征工程
    \item \textbf{Agent 工程}: Agent Loop / Tool Use / Sub-Agent 编排 / Prompt Engineering / MCP 协议 / Memory 系统
    \item \textbf{工程部署}: Git, CLI 工具开发, Vercel Serverless 部署, API 集成 (Anthropic / OpenAI / DeepSeek)
\end{itemize}

%======================================================================
\sectionTitle{工作经历}{\faBriefcase}
%======================================================================
\begin{experiences}
  \experience%
    [2025.12]%
    {2026.05}%
    {广西一三数据科技有限公司 \quad AI 训练师}%
    [\begin{itemize}
      \item 主导自动驾驶 3D 感知数据的标注策略制定，处理叠帧重影与移动物体轨迹连续性，累计处理 100 万+ 目标物，提升环境感知数据质量一致性，支撑感知模型训练
      \item 主导多模态 AIGC 数据评估体系的构建，制定覆盖 15+ 类场景的评估标准，将团队交付合格率提升至 98\%
      \item 负责模型输出的质量分析与人工蒸馏，累计产出高质量偏好数据 100 万+ 条，支撑模型 RLHF 训练与能力迭代
    \end{itemize}]
\end{experiences}

%======================================================================
\sectionTitle{个人项目}{\faRocket}
%======================================================================
\begin{itemize}
    \item[\faCode] \textbf{Claude Code from Scratch — Coding Agent 核心架构复现与扩展}

    从零实现 Claude Code 核心的 Agent Harness 架构，涵盖 Agent Loop、工具系统、上下文压缩与多 Agent 编排等核心模块，总代码量约 5000 行。

    技术栈: TypeScript, Node.js, Anthropic / OpenAI API \\
    GitHub: \href{https://github.com/dok4everak47/claude-code-from-scratch}{\texttt{dok4everak47/claude-code-from-scratch}}
    \begin{itemize}
        \item 完整构建 Agent 执行循环与 13 个内置工具，支持流式并行执行、跨 Provider 消息排序安全网、分层上下文压缩与多类型记忆系统、MCP 协议集成
        \item \textbf{实现 Sub-Agent 分治架构（fork-return）}: 主 Agent 派生 explore / plan / general 三类子 Agent，借工具权限隔离（探索 / 规划仅只读）保护主会话上下文；自定义 agent 发现机制支撑复杂任务自动拆解（核心实现约 200 行）
        \item 开发 Agent Tool System Demo（React + Vite + Tailwind）: 预设场景展示推理与安全网恢复，自由模式接入真实 LLM API 与 Wikipedia / 汇率 / 词典等公开工具，对比模式并行 3 种 System Prompt 并输出量化指标；Vercel Proxy 保护 API Key
        \item 构建 Multi-Agent 可视化演练场（React + TypeScript，SVG 2D 树状编排 + react-three-fiber (Three.js) 3D 场景）: 实时呈现 Sub-Agent 委派 / 进度 / 结果流，内置代码审查、研究报告、客服 3 向分流 3 套协作场景；经 Vercel 部署上线（agent-playground-ruddy.vercel.app）
    \end{itemize}

    \vspace{0.6ex}
    {\color{accentcolor!15}\hrule}
    \vspace{0.6ex}

    \item[\faCode] \textbf{Agent Harness 行为分析与评测工具链}

    自建 Agent 对话日志分析系统，使用 TypeScript + Python（Pandas / NumPy / scikit-learn）处理 5.8 万条对话日志，解决多源日志交错解析、噪声清洗等工程问题。

    技术栈: TypeScript, Python, Pandas, NumPy, scikit-learn
    \begin{itemize}
        \item 基于日志分析提出三层分级记忆模型与动态摘要缓存方案，对应 Harness 中的 Memory 和 Context Engineering 优化方向
        \item 对 Pi / Claude / ChatGPT / DeepSeek 进行对比测试，建立 Agent 行为对比评测框架，产出「情绪-风格解耦对齐」「阈值触发式风格切换」「跨对话情绪」3 项可量化洞察
        \item 分析脚本可扩展用于任意 Agent 对话日志的行为模式提取与量化分析
    \end{itemize}

    \vspace{0.6ex}
    {\color{accentcolor!15}\hrule}
    \vspace{0.6ex}

    \item[\faCode] \textbf{MetroType China — 数据驱动的地铁主题打字应用}

    以中国大陆四城（北京 / 上海 / 广州 / 深圳）25 条线路、684 站的中英文站名为题的打字练习网站，采用完全数据驱动架构，新增城市 / 线路无需改动应用逻辑；用 Python 脚本从公开资料生成结构化站点 / 线路数据，依真实站序与经纬度经 d3-geo (geoMercator) 投影绘制路线，实现城市概览 → 线路详情 → 打字的两级导航与 WPM / CPM 实时反馈。

    技术栈: React 18, Vite 5, pnpm, d3-geo, topojson-client, Python
\end{itemize}

%======================================================================
\sectionTitle{教育背景}{\faGraduationCap}
%======================================================================
\begin{educations}
  \education%
    {2020.09}%
    [2023.06]%
    {桂林电子科技大学}%
    {数学与计算科学学院}%
    {数学}%
    {硕士}
  \separator{0.5ex}
  \education%
    {2016.09}%
    [2020.06]%
    {西安邮电大学}%
    {理学院}%
    {信息与计算科学}%
    {学士}
\end{educations}

\end{document}
